\chapter{Resummation, variational reorganization, and instantons}
\label{chap:resummation-instantons}
\index{resummation|textbf}
\index{divergent series|textbf}
\index{instanton|textbf}

\begin{readercontract}
This appendix is a calculation chapter, not a catalogue of summation
methods.  One benchmark is carried through ordinary truncation, Borel
summation, Pad\'e and Borel--Pad\'e approximation, optimized and variational
perturbation theory, and order-dependent mapping.  The same calculation then
identifies the saddle action which controls large order.  The final section
translates every object into the exact-WKB language used in Stage III.
\end{readercontract}

\calculationroute
{Gaussian integration, analytic continuation, and the Borel transform in
Sec.~\ref{subsec:borel-resummation}.}
{Recover a function from its divergent weak-coupling coefficients, diagnose
the failure of that recovery, and read the missing exponential scale from
large order.}
{The reader can reproduce Figs.~or tables of convergence from the supplied
Wolfram Language file, locate the leading Borel singularity, and explain why
lateral sums and exact-WKB Stokes data are the same kind of information.}

The perturbative series and the exact-WKB series are not rival
approximations.  They are two coordinate descriptions of an analytic
problem.  Ordinary perturbation theory expands a quantity in a coupling;
exact WKB expands a local solution in a semiclassical parameter.  In both
cases the coefficients generally grow factorially, the formal series alone
does not determine a function, and global analytic continuation supplies the
missing information.  The point of the present calculation is to make that
statement operational.

\section{A benchmark with an exact answer}
\label{sec:quartic-zero-dimensional}
\index{quartic integral}
\index{zero-dimensional field theory}

For $\re g>0$, define
\begin{equation}
  I(g)=\frac{1}{\sqrt{2\pi}}
  \int_{-\infty}^{\infty}
  \rme^{-x^2/2-gx^4/4}\dd x .
  \label{eq:quartic-integral}
\end{equation}
This zero-dimensional integral is deliberately elementary.  It has a
Gaussian perturbative vacuum, a nontrivial saddle after analytic
continuation, factorial large order, a Borel singularity, and an exact answer.
It therefore separates the logic of resummation from the numerical and
domain questions of a differential operator.  The repository accompanying
the calculations in this appendix uses the same benchmark
\cite{Morikawa:ResummationRepository}.

The substitution $y=x^2$ and a standard integral representation of the
modified Bessel function give
\begin{equation}
  I(g)=\frac{\rme^{1/(8g)}}{2\sqrt{\pi g}}
  K_{1/4}\!\left(\frac{1}{8g}\right).
  \label{eq:quartic-exact}
\end{equation}
Equation~\eqref{eq:quartic-exact} is a benchmark, not an input to the
resummation algorithms.  In a realistic quantum problem the exact answer is
unknown; one instead compares independent continuations, stability under the
order, and known asymptotic constraints.

\subsection{Derive the weak-coupling coefficients}
\label{subsec:quartic-coefficients}
\index{perturbation series}
\index{Gaussian moment}

Expand only the quartic exponential:
\begin{align}
  I(g)
  &\sim \sum_{n=0}^{\infty}a_n g^n, \\
  a_n
  &=\frac{(-1)^n}{4^n n!\sqrt{2\pi}}
    \int_{-\infty}^{\infty}x^{4n}\rme^{-x^2/2}\dd x \\
  &=(-1)^n\frac{\Gamma(2n+1/2)}{\sqrt{\pi}\,n!}.
  \label{eq:quartic-coefficients}
\end{align}
The second line is a Gaussian moment.  No diagrammatic combinatorics has
been hidden: its rapid growth is the number of Wick contractions in its
simplest form.

Using Stirling's formula in the ratio $a_{n+1}/a_n$ gives
\begin{equation}
  \frac{a_{n+1}}{a_n}
  =-\frac{(2n+1/2)(2n+3/2)}{n+1}
  =-4n\left[1+\order(n^{-1})\right].
  \label{eq:quartic-ratio}
\end{equation}
Thus the coefficients have the schematic behavior
\begin{equation}
  a_n\sim (-4)^n\Gamma(n)\times
  \text{a power of }n.
  \label{eq:quartic-large-order}
\end{equation}
The radius of convergence is zero.  This does not mean that the first few
terms are useless.  It means that truncation is an asymptotic operation whose
order is part of the approximation.

\subsection{Optimal truncation is the first nonperturbative estimate}
\label{subsec:optimal-truncation}
\index{optimal truncation}
\index{least term}

Let $T_n=a_ng^n$.  From Eq.~\eqref{eq:quartic-ratio}, successive terms have
comparable magnitude near
\begin{equation}
  4|g|n\simeq 1,
  \qquad
  n_\star\simeq \frac{1}{4|g|}.
  \label{eq:least-term-order}
\end{equation}
Truncating substantially beyond $n_\star$ makes the approximation worse.
The size of the least term is exponentially small,
\begin{equation}
  |T_{n_\star}|\sim
  \rme^{-1/(4|g|)}\times \text{a power of }g.
  \label{eq:least-term-scale}
\end{equation}
Already at this stage perturbation theory predicts the scale of the effect
which it cannot resolve.  Section~\ref{sec:instanton-large-order} will derive
the same number $1/4$ from a nontrivial saddle.

\begin{checkpointbox}
A divergent series should never be judged from one fixed truncation order.
Plot the partial sums against the order.  The useful region is the plateau
near the least term; the later divergence is part of the prediction, not a
software failure.
\end{checkpointbox}

\section{Borel and lateral summation}
\label{sec:borel-lateral-appendix}
\index{Borel transform|textbf}
\index{Borel sum|textbf}
\index{lateral sum}

For a formal series $F(g)\sim\sum_{n\geq0}a_ng^n$, define the Borel transform
and a directional Laplace transform by
\begin{align}
  \widehat F(t)
  &=\sum_{n=0}^{\infty}\frac{a_n}{n!}t^n,
  \label{eq:appendix-borel-transform}\\
  \mathcal S_\theta F(g)
  &=\frac{1}{g}\int_{0}^{\rme^{\rmi\theta}\infty}
    \rme^{-t/g}\widehat F(t)\dd t.
  \label{eq:appendix-borel-laplace}
\end{align}
Termwise integration recovers the original formal series because
\begin{equation}
  \frac{1}{g}\int_0^\infty \rme^{-t/g}t^n\dd t=n!g^n.
  \label{eq:borel-inverse-check}
\end{equation}
The factorial division in Eq.~\eqref{eq:appendix-borel-transform} repairs the
leading divergence, while the Laplace integral restores the coefficients.

For the quartic integral, Eq.~\eqref{eq:quartic-large-order} implies that the
nearest singularity of $\widehat I(t)$ lies at
\begin{equation}
  t=-\frac14.
  \label{eq:quartic-borel-singularity}
\end{equation}
For positive real $g$, this singularity is away from the positive Laplace
ray, and the perturbative series is Borel summable.  For negative real $g$,
the relevant ray meets the continued singularity.  One must specify whether
the contour passes above or below it:
\begin{equation}
  \mathcal S_{0^\pm}F(g)
  =\lim_{\epsilon\downarrow0}
  \frac{1}{g}\int_{0}^{\rme^{\pm\rmi\epsilon}\infty}
  \rme^{-t/g}\widehat F(t)\dd t.
  \label{eq:lateral-borel-sums}
\end{equation}
Their difference is exponentially small in $g$.  The two values are not a
defect of the formalism: they are the analytic record of choosing a side of a
Stokes ray.  A complete physical problem supplies compensating saddle or
boundary-condition data so that the final observable has the required
reality or outgoing property; this is the elementary form of a transseries
completion \cite{Dorigoni:2014hea,Delabaere:1999xx}.

\subsection{What the Borel plane tells us before integration}
\label{subsec:borel-diagnostics}
\index{Borel plane}
\index{Borel singularity}
\index{large-order relation}

The Borel plane is a map of competing saddles.  If, locally,
\begin{equation}
  \widehat F(t)\sim \frac{C}{(1-t/A)^\beta},
  \label{eq:borel-local-singularity}
\end{equation}
then singularity analysis gives
\begin{equation}
  a_n\sim
  \frac{C}{\Gamma(\beta)}A^{-n}\Gamma(n+\beta).
  \label{eq:borel-large-order-dictionary}
\end{equation}
The position $A$ determines the exponential $\rme^{-A/g}$; the exponent
$\beta$ determines the power-law prefactor; the discontinuity coefficient
determines the Stokes multiplier.  Numerical large-order data can therefore
be used in the reverse direction: coefficient ratios estimate $A$, and
subleading Richardson-type fits estimate $\beta$ and corrections.

\begin{warningbox}
A numerical Laplace integral through a Pad\'e pole is not a definition of a
Borel sum.  First inspect the poles, then choose a direction or a specified
principal-value prescription.  An integration routine which silently steps
around a pole has made a physical choice on the user's behalf.
\end{warningbox}

\section{Pad\'e and Borel--Pad\'e approximation}
\label{sec:pade-borel-pade}
\index{Pad\'e approximant|textbf}
\index{Borel--Pad\'e approximation|textbf}

Given $N+1$ coefficients, the $[L/M]$ Pad\'e approximant is the rational
function
\begin{equation}
  [L/M]_F(g)=\frac{P_L(g)}{Q_M(g)},
  \qquad L+M=N,
  \qquad Q_M(0)=1,
  \label{eq:pade-definition}
\end{equation}
whose Taylor expansion agrees with $F$ through order $N$.  Pad\'e
approximation can represent poles and mimic a branch cut by a sequence of
zeros and poles.  It does not, by itself, remove factorial divergence.

The robust combination is therefore:
\begin{enumerate}[label=\arabic*.]
  \item divide $a_n$ by $n!$ to form a truncated Borel transform;
  \item replace that polynomial by a near-diagonal Pad\'e approximant;
  \item inspect its pole pattern in the Borel plane;
  \item Laplace integrate along a prescribed contour.
\end{enumerate}
Explicitly,
\begin{align}
  \widehat F_N(t)
  &=\sum_{n=0}^{N}\frac{a_n}{n!}t^n,
  \\
  \widehat F_N^{[L/M]}(t)
  &= [L/M]_{\widehat F_N}(t),
  \\
  \vsup{F_N}{BP}(g)
  &=\frac{1}{g}\int_0^\infty
    \rme^{-t/g}\widehat F_N^{[L/M]}(t)\dd t.
  \label{eq:borel-pade-algorithm}
\end{align}
For Eq.~\eqref{eq:quartic-integral}, stable near-diagonal sequences recover
Eq.~\eqref{eq:quartic-exact} for positive $g$.  Alternating poles and zeros
accumulate along the negative Borel axis, resolving the branch cut which
starts at Eq.~\eqref{eq:quartic-borel-singularity}.

\subsection{A convergence protocol rather than one decimal number}
\label{subsec:resummation-protocol}
\index{convergence!resummation}
\index{spurious pole}

For each $g$, calculate a family of $[L/M]$ approximants with
$|L-M|\leq1$.  Record the result together with the closest pole to the
integration ray.  A trustworthy digit is one which survives:
\begin{itemize}
  \item increasing $N$;
  \item exchanging $[L/M]$ and $[M/L]$;
  \item a small contour rotation which crosses no singularity;
  \item increased working precision;
  \item comparison with an independent method such as ODM.
\end{itemize}
Near-cancelling pole--zero pairs are often finite-order artifacts.  They may
be diagnosed, but should not be deleted without reporting the criterion.

\section{Optimized and variational perturbation theory}
\label{sec:opt-vpt}
\index{optimized perturbation theory|textbf}
\index{variational perturbation theory|textbf}
\index{delta expansion}
\index{principle of minimal sensitivity}

Resummation can also be organized by changing the unperturbed problem.  Add
and subtract a solvable term containing a trial parameter $\Omega$, insert a
bookkeeping parameter $\delta$, expand at fixed $\Omega$, set $\delta=1$,
and only then choose $\Omega$.  For the quartic integral, write
\begin{equation}
  \frac{x^2}{2}+\frac{gx^4}{4}
  =\frac{\Omega x^2}{2}
  +\delta\left[
    \frac{(1-\Omega)x^2}{2}+\frac{gx^4}{4}
  \right].
  \label{eq:delta-splitting}
\end{equation}
At order $N$, expand the second exponential through $\delta^N$ and evaluate
all moments with the Gaussian of width $\Omega^{-1/2}$.  Denote the result
after setting $\delta=1$ by $I_N(g,\Omega)$.

The exact integral is independent of $\Omega$, but the truncated result is
not.  The principle of minimal sensitivity chooses a locally stationary
approximant,
\begin{equation}
  \frac{\partial I_N(g,\Omega)}{\partial\Omega}=0,
  \label{eq:pms-condition}
\end{equation}
while fastest apparent convergence sets a selected highest-order correction
to zero.  Multiple stationary points are common.  One follows a continuous
root sequence with $N$ and reports the selection rule; choosing the most
favorable root after seeing the exact answer is not an algorithm.

Optimized perturbation theory emphasizes the invariance of a complete answer
under the artificial parameter and uses stationarity as the finite-order
criterion \cite{Stevenson:1981vj}.  The linear delta expansion makes the
bookkeeping construction explicit, with convergence results available in
important model problems \cite{Buckley:1992pc,Duncan:1992ba}.  Variational
perturbation theory further builds known strong-coupling scaling into the
reorganization \cite{Kleinert:1995hc}.  The names overlap in the literature;
the defining data are therefore the splitting, the truncation rule, and the
optimization condition, not the label.

\subsection{Why optimization can recover strong coupling}
\label{subsec:vpt-strong-coupling}
\index{strong-coupling expansion}

Rescale $x=g^{-1/4}y$ in Eq.~\eqref{eq:quartic-integral}.  Then
\begin{equation}
  I(g)=g^{-1/4}\frac{1}{\sqrt{2\pi}}
  \int_{-\infty}^{\infty}
  \rme^{-y^4/4-y^2/(2\sqrt g)}\dd y,
  \label{eq:quartic-strong-coupling}
\end{equation}
so the leading strong-coupling power is $g^{-1/4}$ and corrections occur in
powers of $g^{-1/2}$.  A variational mapping which respects these exponents
has information unavailable to a raw weak-coupling polynomial.  In operator
problems the analogous information comes from dimensional analysis,
semiclassical scaling, or a solvable strong-coupling limit.

\begin{checkpointbox}
Optimization does not make an approximation variational in the sense of an
upper bound.  Unless a Rayleigh--Ritz theorem applies to the precise
functional being optimized, stationarity supplies stability, not an
inequality.
\end{checkpointbox}

\section{Order-dependent mapping}
\label{sec:odm-appendix}
\index{order-dependent mapping|textbf}
\index{conformal map}

Order-dependent mapping combines analytic continuation with an optimized
parameter.  Introduce
\begin{equation}
  g=\rho\frac{\lambda}{(1-\lambda)^\alpha}.
  \label{eq:odm-general-map}
\end{equation}
The exponent $\alpha$ is chosen from the analytic or strong-coupling
behavior.  For the present quartic integral a useful map is
\begin{equation}
  g=\rho\frac{\lambda}{(1-\lambda)^2}.
  \label{eq:odm-quartic-map}
\end{equation}
Substitute Eq.~\eqref{eq:odm-quartic-map} into the weak-coupling series,
reexpand in $\lambda$ through order $N$, and write
\begin{equation}
  I_N(g;\rho)=\sum_{n=0}^{N}c_n(\rho)\lambda(g,\rho)^n.
  \label{eq:odm-truncated}
\end{equation}
The parameter is allowed to depend on the order.  A common condition is
\begin{equation}
  c_N(\rho_N)=0,
  \label{eq:odm-fac}
\end{equation}
or, alternatively, stationarity with respect to $\rho$.  As $N$ grows,
$\rho_N$ moves so that the finite analytic domain in $\lambda$ resolves a
larger part of the cut $g$ plane.  The general construction and its
convergence mechanisms are reviewed in Ref.~\cite{Zinn-Justin:2010zzb}; the
specific implementation used here follows
Ref.~\cite{Morikawa:ResummationRepository}.

\subsection{How ODM differs from a fixed conformal Borel map}
\label{subsec:odm-versus-conformal}
\index{conformal Borel mapping}

A conformal Borel method maps a known cut Borel domain to a disk using a
fixed singularity location.  ODM maps the coupling itself and changes the
scale $\rho_N$ with the order.  The former is strongest when the Borel
singularity geometry is known; the latter can infer an efficient scale from
the coefficient sequence and strong-coupling exponent.  They may be
combined, but then each source of input must be stated to avoid counting the
same analytic information twice.

\section{Instantons and the origin of large order}
\label{sec:instanton-large-order}
\index{instanton action|textbf}
\index{saddle point}
\index{large-order behavior}

Continue the quartic coupling to $g=-|g|$.  The exponent can be written as
\begin{equation}
  S(x)=\frac{x^2}{2}-\frac{|g|x^4}{4}.
  \label{eq:continued-quartic-action}
\end{equation}
Besides $x=0$, the stationary points satisfy
\begin{equation}
  S'(x)=x-|g|x^3=0,
  \qquad
  x_\star=\pm\frac{1}{\sqrt{|g|}}.
  \label{eq:quartic-saddles}
\end{equation}
Their action relative to the perturbative saddle is
\begin{equation}
  S(x_\star)-S(0)=\frac{1}{4|g|}.
  \label{eq:quartic-saddle-action}
\end{equation}
This is precisely the exponential scale inferred from the least term in
Eq.~\eqref{eq:least-term-scale}, and its reduced action $A=1/4$ is precisely
the Borel singularity in Eq.~\eqref{eq:quartic-borel-singularity}.  Three
calculations have found the same invariant:
\begin{equation}
  \text{large-order ratio}
  \quad\longleftrightarrow\quad
  \text{Borel singularity}
  \quad\longleftrightarrow\quad
  \text{saddle action}.
  \label{eq:large-order-triangle}
\end{equation}

In quantum mechanics the nontrivial saddle is a finite-action trajectory in
Euclidean time rather than a point.  Expanding about it gives a fluctuation
determinant and zero-mode measure.  Multi-instanton sectors generate powers
of $\rme^{-A/g}$, possible logarithms, and interactions between saddles
\cite{Zinn-Justin:1981qzi,Marino:2015Instantons}.  The perturbative coefficients around one saddle
encode fluctuations around the others.  This is the practical content of
resurgence, not merely the statement that a divergent series can sometimes
be summed.

\subsection{Dispersion relation for the coefficients}
\label{subsec:dispersion-large-order}
\index{dispersion relation!large order}

Suppose $F(g)$ is analytic in a cut plane and has a discontinuity on the
negative axis.  Cauchy's formula gives, under the appropriate subtractions,
\begin{equation}
  a_n
  =\frac{(-1)^{n+1}}{\pi}
  \int_0^\infty
  \frac{\im F(-s+\rmi0)}{s^{n+1}}\dd s.
  \label{eq:large-order-dispersion}
\end{equation}
If the discontinuity has the semiclassical form
\begin{equation}
  \im F(-s+\rmi0)
  \sim C s^{-b}\rme^{-A/s}
  \left[1+\order(s)\right],
  \label{eq:instanton-discontinuity}
\end{equation}
then $u=A/s$ in Eq.~\eqref{eq:large-order-dispersion} yields
\begin{equation}
  a_n\sim
  \frac{(-1)^{n+1}C}{\pi}
  A^{-n-b}\Gamma(n+b)
  \left[1+\order(n^{-1})\right].
  \label{eq:instanton-large-order}
\end{equation}
Thus the imaginary part of an unstable continuation controls the late terms
of the stable perturbative problem.  For a resonance, the same imaginary
part becomes a decay width after the outgoing boundary value has been fixed.

\section{The exact-WKB translation}
\label{sec:resummation-wkb-dictionary}
\index{exact WKB!resummation dictionary}
\index{Voros period}
\index{Stokes automorphism}
\index{transseries}

The exact-WKB construction in Sec.~\ref{sec:exact-wkb} begins from a local
Riccati series, Borel resums it in a direction, and transports the resulting
basis through Stokes sectors.  The global connection coefficient
$\Cincoming(E)$ is then an analytic function assembled from local data.  The
dictionary is exact at the structural level:
\begin{center}
\begin{tabularx}{\textwidth}{@{}>{\raggedright\arraybackslash}p{0.45\textwidth}
  >{\raggedright\arraybackslash}X@{}}
\toprule
Coupling expansion & Exact-WKB connection problem \\
\midrule
formal perturbative coefficients & formal Riccati coefficients \\
\addlinespace
Borel singularity at action $A$ & finite action path or WKB cycle \\
\addlinespace
lateral Borel sum & directional exact-WKB sum $\mathcal S_{\theta^\pm}$ \\
\addlinespace
Stokes discontinuity & change of the resummed WKB basis \\
\addlinespace
instanton sector & exponential of a Voros period $\VorosExp_\gamma$ \\
\addlinespace
physical transseries parameter & boundary or connection data \\
\addlinespace
zero of a continued denominator & zero of $\Cincoming(E)$ \\
\bottomrule
\end{tabularx}
\end{center}

The last two rows are the reason this appendix belongs in the present book.
For a generic perturbative observable, a transseries parameter can look like
an extra constant.  In a scattering problem it is fixed by preparation and
asymptotic boundary conditions.  The outgoing condition selects a lateral
continuation and imposes
\begin{equation}
  \Cincoming(\resonantE)=0.
  \label{eq:appendix-resonance-condition}
\end{equation}
The resonance is therefore not obtained by adding a complex energy by hand.
It is the zero of the globally continued coefficient.  Its residue depends
on $\partial_E\Cincoming$, so the fluctuation and normalization information
which accompanies the exponential sector remains observable.

\subsection{What ordinary resummation cannot decide}
\label{subsec:resummation-limit}
\index{boundary condition!and resummation}

A coefficient list does not specify:
\begin{itemize}
  \item the physical sheet of the energy or momentum plane;
  \item which asymptotic channel is incoming or outgoing;
  \item the test space on which a continued eigenfunctional acts;
  \item whether a Borel ambiguity is cancelled, retained as a width, or
        moved by a change of observable;
  \item which continuum background accompanies a pole.
\end{itemize}
These are connection, domain, and representation data.  Exact WKB is useful
because it keeps them in the calculation from the beginning.

\section{A reproducible Wolfram Language laboratory}
\label{sec:wolfram-resummation-lab}
\index{Wolfram Language}
\index{reproducibility!resummation}

The file \texttt{notebooks/resummation\_workflow.wl} supplied with the source
archive implements the coefficient generation, direct truncation,
Borel--Pad\'e summation, Borel-plane pole inspection, and the ODM map.  It is
a compact, text-readable companion to the larger public notebook collection
\cite{Morikawa:ResummationRepository}.  A separate Wolfram Community article
demonstrates how symbolic and numerical code can accompany the exact-WKB
resonance analysis itself \cite{Morikawa:2026WolframCommunity}.

The core commands are included here so that the printed book remains
self-contained:
\lstinputlisting[
  basicstyle=\ttfamily\footnotesize,
  breaklines=true,
  columns=fullflexible,
  frame=single,
  caption={Minimal Wolfram Language workflow for the quartic integral.},
  label={lst:resummation-workflow}
]{notebooks/resummation_workflow.wl}

\subsection{Protocol for comparing methods}
\label{subsec:method-comparison-protocol}
\index{method comparison}

Use exactly the same coefficient list and working precision for every method.
For $g=0.2$ and $g=0.4$, carry out the following steps.
\begin{enumerate}[label=\arabic*.]
  \item Evaluate Eq.~\eqref{eq:quartic-exact} at high precision only to create
        withheld validation data.
  \item Plot partial sums for all $0\leq N\leq N_{\max}$ and mark the least
        term predicted by Eq.~\eqref{eq:least-term-order}.
  \item Compute near-diagonal Borel--Pad\'e sequences and reject no pole
        without recording its pole--zero distance.
  \item Follow one continuous PMS or FAC root for OPT/VPT as $N$ changes.
  \item Follow one continuous positive ODM root $\rho_N$ and record both
        $\rho_N$ and the mapped variable $\lambda(g,\rho_N)$.
  \item Compare errors against the exact answer only after all selection
        rules have been fixed.
\end{enumerate}
Representative exact values are
\begin{align}
  I(0.2)&=0.9079847774\ldots,\\
  I(0.4)&=0.8576085854\ldots.
  \label{eq:quartic-benchmark-values}
\end{align}
The purpose is not to declare a universal winner.  Borel--Pad\'e is strongest
when the Borel geometry is well resolved; VPT is strongest when reliable
strong-coupling exponents are available; ODM is attractive when a useful map
and a stable order-dependent root can be identified.  Exact WKB adds the
global channel geometry needed for spectral and scattering questions.

\begin{derivationbox}{Exercises}
\begin{enumerate}[label=\textbf{R.\arabic*}]
  \item Derive Eq.~\eqref{eq:quartic-exact} from
        Eq.~\eqref{eq:quartic-integral} and verify its $g\downarrow0$
        asymptotic expansion.
  \item Use Eq.~\eqref{eq:quartic-coefficients} to estimate the Borel
        singularity from three successive coefficient ratios.  Determine the
        leading $1/n$ correction.
  \item Compare direct Pad\'e and Borel--Pad\'e at $g=0.4$.  Plot all Borel
        poles and explain which sequences can be integrated on the positive
        axis without a prescription.
  \item Derive $I_N(g,\Omega)$ through $N=3$ from
        Eq.~\eqref{eq:delta-splitting}.  Find all real positive PMS roots and
        track them from $g=0.05$ to $g=1$.
  \item Implement Eqs.~\eqref{eq:odm-quartic-map}--\eqref{eq:odm-fac} without
        using the supplied function.  Explain why selecting a different root
        independently at each order destroys the meaning of convergence.
  \item Starting from Eq.~\eqref{eq:large-order-dispersion}, derive
        Eq.~\eqref{eq:instanton-large-order}, including the overall power of
        $A$.
  \item For the inverted Rosen--Morse problem in
        Sec.~\ref{subsec:rosen-morse}, identify the WKB action which replaces
        $A=1/4$ and explain how its Stokes multiplier enters
        $\Cincoming(E)$.
\end{enumerate}
\end{derivationbox}

\section{What to retain for the rest of the book}
\label{sec:resummation-retain}
\index{resummation!summary}

The reusable lesson is a five-step discipline.
\begin{enumerate}[label=\arabic*.]
  \item Determine the domain, coupling variable, and asymptotic boundary
        condition before manipulating coefficients.
  \item Diagnose late terms and locate candidate action scales.
  \item Choose a summation or reorganization whose analytic assumptions are
        explicit.
  \item Test order, precision, contour, root, and representation stability.
  \item Translate the recovered analytic function back into invariant pole,
        residue, and continuum data.
\end{enumerate}
This discipline is the perturbative analogue of the calculation spine in
Sec.~\ref{sec:computational-spine}: local data are useful only after a global
connection prescription has been fixed.
